\documentclass[11pt]{article}
\usepackage[letterpaper,margin=1in]{geometry}
\usepackage{amsmath,amssymb,amsthm,microtype}
\usepackage[hidelinks]{hyperref}
\hypersetup{pdftitle={Initial omnific models of discrete initial surreal groups},pdfauthor={Alec Kriebel},pdfsubject={An answer to the discrete initial group question of Kaplan and Ehrlich},pdfkeywords={surreal numbers, initial subgroup, omnific integers, Hahn series}}
\newcommand{\No}{\mathbf{No}}
\newcommand{\Oz}{\mathbf{Oz}}
\newcommand{\D}{\mathbb{Z}[1/2]}
\newcommand{\supp}{\operatorname{supp}}
\newcommand{\Z}{\mathbb{Z}}
\newcommand{\R}{\mathbb{R}}
\newtheorem{theorem}{Theorem}
\newtheorem{lemma}[theorem]{Lemma}
\newtheorem{proposition}[theorem]{Proposition}
\title{Initial omnific models of discrete initial surreal groups}
\author{Alec Kriebel\\\small Independent researcher\\\small ORCID: \href{https://orcid.org/0009-0001-9320-500X}{0009-0001-9320-500X}}
\date{17 October 2026}
\begin{document}
\maketitle
\begin{abstract}
Every discrete initial ordered additive subgroup of the surreal numbers is
order-group-isomorphic to a subgroup of the omnific integers that is initial
in the surreal numbers. This holds for sets and proper classes in NBG with
global choice and answers the discrete-group question posed by Kaplan and
Ehrlich. An explicit operation on ordinal-length sign sequences moves the
least exponent to zero without creating a right-simplicity-predecessor
obligation. Relabeling the actual normal-form supports and normalizing
only the bottom coefficient group then gives the required image.
An appendix supplies a detailed sufficiency proof of the initial-group
characterization used in the argument, addressing support exhaustion and
the legality of limit-cut options.
\end{abstract}

\section{Question, conventions and established inputs}
An \emph{initial} subclass of $\No$ contains every proper sign-sequence
prefix of each of its members. This is closure under simplicity, not
numerical convexity. We write $u\prec_s v$ for a proper prefix,
$u\preceq_s v$ when equality is also allowed,
$P(v)=\{u:u\prec_s v\}$, and
$R_s(v)=\{u\in P(v):u>v\}$. All orders and group operations are inherited
from $\No$. A nontrivial ordered group is \emph{discrete} if it has a least
positive element; the zero group is treated separately.

The following is the full additive-group question, with no multiplicative
closure hypothesis.
\begin{theorem}\label{main}
Let $G\subseteq\No$ be a discrete initial ordered additive subgroup.
There is an order-group isomorphism $T:G\longrightarrow H$, where
$H\subseteq\Oz$ is initial in $\No$. The assertion holds when $G$ is a set
and when it is a proper class. The zero group also has the stated property.
\end{theorem}
The target occurs in Kaplan's April 2015 honors thesis
\cite[\S7.1, Question 1, p.~47]{thesis}, in his joint-work contribution
with Ehrlich \cite[Question 1, p.~3356]{owr}, and in Ehrlich--Kaplan
\cite[Question 9.1 in the author version]{ek}.
The numbering in arXiv:1512.04001v1 is Question 2.
The isomorphism in Theorem~\ref{main} is not required to preserve simplicity.

We work in NBG with global choice, as in \cite{ek}. Every surreal has a
unique normal form
\[
 x=\sum_{\mu<\lambda}r_\mu\omega^{y_\mu},
 \qquad r_\mu\in\R\setminus\{0\},\quad y_\mu\text{ strictly decreasing},
\]
with set ordinal $\lambda$; zero has empty support. Addition is
coefficientwise and order is determined by the leading nonzero coefficient.
All coefficient products below are real scalings of monomials.
We identify groups with their \emph{canonical} normal-form images,
not with an arbitrarily chosen abstract Hahn representation.
We use these established results of \cite{ek}:
\begin{enumerate}
\item The canonical normal-form correspondence is an ordered-group
isomorphism (Proposition 3.3).
\item A nonzero initial subgroup of $\R$ is either $2^{-n}\Z$ for
$n\in\mathbb{N}_0$, or contains $\D$ (Lemma 5.1).
\item A nonzero additive subgroup $A$ of normal forms is initial precisely
when it is truncation closed and cross sectional over an initial exponent
class $\Gamma$, its pure coefficient groups
\[
 R_y=\{r\in\R:r\omega^y\in A\}
\]
are initial in $\R$, and $\D\subseteq R_y$ whenever
$x,y\in\Gamma$ and $y\in R_s(x)$ (Theorem 5.1).
Here cross sectional means $\omega^y\in A$ for every $y\in\Gamma$.
\item A normal form lies in $\Oz$ precisely when its exponents are
nonnegative and its coefficient at zero, if present, is integral
(\S6.2).
\end{enumerate}
For truncation-closed groups, every coefficient occurring in an actual
series belongs to its pure coefficient group: subtract its two adjacent
truncations, using the full series as the later truncation at a last term.
The exponent class is therefore exactly the class of occurring exponents
when the group is also cross sectional.
The appendix checks the sufficiency direction in item 3 without
a bounded-support assumption or fractional options outside the group.
The underlying standard normal-form and monomial-cut facts remain
explicit inputs.

\section{An exponent map with the required prefix control}
\begin{lemma}\label{bottom}
If $G$ is nontrivial, discrete and initial, its exponent class $\Gamma$
has minimum $m=-\alpha$ for a set ordinal $\alpha$, and
$R_m=2^{-n}\Z$ for some $n\in\mathbb{N}_0$.
Its least positive element is $2^{-n}\omega^m$.
\end{lemma}
\begin{proof}
Let $g>0$ be least, with leading term $r\omega^m$, $r>0$.
Truncation closure puts this term and the residual tail in $G$.
A nonzero tail has absolute value strictly between zero and $g$,
a contradiction. Thus $g=r\omega^m$.
The group $R_m$ has least positive element $r$, so item 2 gives
$R_m=2^{-n}\Z$ and $r=2^{-n}$.
If $y\in\Gamma$ were below $m$, cross sectionality would give
$0<\omega^y<g$. Hence $m=\min\Gamma$.
A plus sign in $m$ would make the prefix before that sign a member of
$\Gamma$ strictly below $m$. Thus $m$ is the all-minus sequence
$(-)^\alpha$, namely $-\alpha$.
\end{proof}
The least-positive-element form is already present in
\cite[Proposition 5.2]{ek}; the negative-ordinal minimum observation
also appears in \cite[Lemma 3.2]{adh}. The proof records their
precise role here.

Concatenation of sign sequences is denoted by juxtaposition, and $e$
is the empty sequence. Concatenation is the ordered sum of blocks;
no cancellation of ordinal lengths is intended.
\begin{lemma}\label{rotation}
Let $\Gamma$ be initial with minimum $a=(-)^\alpha$.
There is a strictly increasing map $f:\Gamma\to\No$ such that
$\Delta=f[\Gamma]$ is initial, $\min\Delta=0$, and
\begin{equation}\label{right}
 f(y)\in R_s(f(x))\quad\Longrightarrow\quad y\in R_s(x).
\end{equation}
\end{lemma}
\begin{proof}
In fact the formula is defined on the full ray $s\ge a$:
\begin{equation}\label{formula}
 f(s)=
 \begin{cases}
 e,&s=a,\\
 (+)s,&s>a\text{ and }a\text{ is not a prefix of }s,\\
 (+)at,&s=a(+)t.
 \end{cases}
\end{equation}
The cases exhaust the ray because an extension of $a$ above $a$ has
a plus as its next sign. Call the second case A and the third B.

Within each case, replacement of a common prefix preserves and reflects
order and proper-prefix relations. For a case-A sequence $s$ and a
case-B sequence $u$, either $s=(-)^\beta$ with $\beta<\alpha$, or
$s$ first deviates from $a$ by a plus at a position $\beta<\alpha$.
In the former case $s>u$ and $s\prec_s u$; its image is likewise a
proper prefix above $f(u)$, whose next sign is minus. In the latter
case $s>u$, with no prefix comparability, and the same first-deviation
comparison holds after the common prepended plus. These alternatives
also rule out collisions between the two cases.
Thus $f$ is strictly increasing and, away from $a$, preserves and reflects
proper prefixes. Its value at $a$ is zero and all other images are positive.

For every $s>a$ we have the exact identity
\begin{equation}\label{prefix}
 P(f(s))=\{e\}\ \cup\ \{f(v):v\in P(s),\ v\ne a\}.
\end{equation}
For case A, strip the initial plus. For case B, the proper prefixes of
$(+)at$ are exactly $e$, the sequences $(+)(-)^\beta$ with
$\beta<\alpha$, and $(+)av$ with $v$ a proper prefix of $t$.
They are respectively $f(a)$, $f((-)^\beta)$, and $f(a(+)v)$.
When $t=e$, the last family is empty: $(+)a$ is the complete image,
not an additional proper prefix.
Equation~\eqref{prefix} and initiality of $\Gamma$ give initiality of
$\Delta$. Finally, a right predecessor in $\Delta$ cannot be zero,
since zero is its minimum. Prefix reflection away from $a$ and
strict monotonicity give \eqref{right}.
\end{proof}
The argument treats arbitrary ordinal lengths. For instance
$1+\alpha=\alpha$ can hold at a limit, but the ordered-block prefix
partition in the proof is unchanged. The full-ray inverse is explicit:
$h(e)=a$, and, for $z=(+)v$, set $h(z)=a(+)t$ if $v=at$,
and $h(z)=v$ otherwise. The two compositions are the identity.
For $\alpha=0$, \eqref{formula} is the identity on nonnegative surreals.
No completion of $\Gamma$ or transfinite class recursion is used.

\section{Transport of the actual subgroup}
\begin{proof}[Proof of Theorem~\ref{main}]
The zero group uses the identity. For a nonzero group take $m,n,f$ from
Lemmas~\ref{bottom} and~\ref{rotation}. Set $c_m=2^n$ and $c_y=1$ for $y\ne m$.
Define
\begin{equation}\label{transport}
 T\left(\sum_{\mu<\lambda}r_\mu\omega^{y_\mu}\right)
    =\sum_{\mu<\lambda}c_{y_\mu}r_\mu\omega^{f(y_\mu)}.
\end{equation}
The support on the right is a set, with the same reverse well-order type
as the original. Thus it is a normal form of a surreal.
Positive fixed scalings and injective increasing exponent relabeling make
$T$ additive, injective and order preserving. Set $H=T[G]$.
This is exactly the image of the actual group; no full Hahn group is added.

Its exponent class is $\Delta$, and truncation closure follows because
corresponding truncations commute with $T$.
For $y\ne m$, $T(\omega^y)=\omega^{f(y)}$; at the bottom,
$T(2^{-n}\omega^m)=1$. Hence $H$ is cross sectional.
A one-term image has a one-term preimage, so its pure coefficient groups
are exactly
\begin{equation}\label{coeff}
 R'_0=\Z,\qquad R'_{f(y)}=R_y\quad(y\ne m).
\end{equation}
They are initial. If $f(y)\in R_s(f(x))$, then $y\in R_s(x)$ by
\eqref{right}, and $y\ne m$. Thus $\D\subseteq R_y=R'_{f(y)}$.
Every hypothesis of the initial-group characterization holds, and $H$
is initial in $\No$. Its exponents are nonnegative. By truncation
closure, every occurring zero coefficient belongs to $R'_0=\Z$.
The omnific criterion therefore gives $H\subseteq\Oz$.

For a proper class $G$, the parameters $\alpha,n$ are still sets.
The explicit formula \eqref{formula} defines a class function, and
class replacement sends each individual support to a set.
Equation~\eqref{transport} defines $T$ and its image class by elementary
comprehension. Every group calculation uses finitely many set supports.
The characterization applies to classes as justified in the appendix.
There is no proper-class support, class-sized sum, or stronger
class-recursion assumption.
\end{proof}

\section{Attribution, verification and limits}
Exponent relabeling and positive coefficient scaling are established
Hahn-group mechanisms; an explicit later account is
\cite[Proposition 3.9 and Construction 3.10]{kks}.
General order/simplicity substructures and partition constructions
\cite{bvh} do not by themselves guarantee the additive compatibility and
initial-image property needed here. The specific bridge proved above is
the sign operation \eqref{formula}, its right-prefix control, and its use
in the actual subgroup. No claim of first discovery of the general
transport machinery is made.

The source package records extensive bounded primary-source priority
searches by the author and three internal AI review approaches.
No earlier explicitly recorded full solution was located in the inspected
accessible material. This is not a certificate of novelty, firstness,
causal independence or continued global openness.
Some publisher versions were not authenticated against the author/arXiv
versions. Material unread fulltext leads include Ehrlich's 2024 handbook
chapter \cite{handbook}, the full 2001 predecessor and other access-limited
works specified in the priority ledger. Only the 2024 chapter's primary
abstract and references were read; no absence claim is made about its
unread contents.

The portable verification program checks finite sign trees, prefix
identities, right-predecessor implications, rational finite-support
additivity and bottom normalization. Ordinary and optimized runs compare
the full expected output. Four altered constructions fail computed
properties in each mode. These finite diagnostics do not establish
the ordinal-length or proper-class theorem, which is proved above.
Historical reports mention unavailable earlier programs; those are not
represented as reproduced by this package.

AI tools were used extensively in developing, writing, researching,
reproducing and adversarially checking this work. Three initial full
mathematical approaches and a further focused dependency reconstruction
were internally reviewed; exact exposure histories and source-reading
limits are recorded in the package. This preprint is unrefereed.
Internal AI reviews are not external human peer review, journal
acceptance or formal proof-assistant certification.

\appendix
\section{A detailed proof of the initial-group sufficiency criterion}
\label{appendix}
This appendix proves the sufficiency assertion used above. The
characterization is credited to \cite[Theorem 5.1]{ek}, not claimed as new.
In the inspected author proof two literal steps require care:
the union of bounded exponent stages need not exhaust cofinal supports,
and dyadic perturbations of a series need not belong to the given
coefficient groups. The proof below avoids both steps. This concerns
proof presentation, not a counterexample to the theorem.

\begin{proposition}\label{suff}
Let $A$ be a truncation-closed, cross-sectional additive subgroup
of normal forms over an initial exponent class $\Gamma$.
Suppose every $R_y$ is initial in $\R$ and
$\D\subseteq R_y$ for $y\in R_s(x)$, $x,y\in\Gamma$.
Then the corresponding subgroup of $\No$ is initial.
\end{proposition}
We use standard recursive addition and negation, the simplicity theorem,
normal-form uniqueness and comparison, the monomial cut
\begin{equation}\label{leadercut}
 \omega^y=\{\,0,\ k\omega^{y^L}\mid 2^{-k}\omega^{y^R}\,\},
\end{equation}
where $k=1,2,\ldots$, $y^L\in P(y)$ with $y^L<y$, and $y^R\in R_s(y)$,
the pure-multiple cuts of \cite[Lemma 3.1(ii)--(iii)]{ek}, and the usual
infinite limit-series cut of \cite[Proposition 3.2(iii)]{ek}.
These normal-form facts, and the real-coefficient classification, are
imports; this appendix does not rederive all foundational surreal theory.

\subsection{Initial sums and adjoining a cut}
If $q=\{L\mid R\}$ and $z\prec_s q$, then $z$ is a prefix of or equal to
an option on its side. Indeed, if $z<q$ and every $l\in L$ were below
$z$, then $L<z<R$, contrary to simplicity of $q$. Thus
$z\le l<q$ for some $l\in L$. A sign sequence lying between a sequence
and one of its extensions must extend the shorter sequence, by the
first differing sign comparison. Consequently $z\preceq_s l$.
For $z>q$ use a member $r\in R$ with $q<r\le z$.
In particular, if all options belong to an initial set $M$, all proper
prefixes of $q$ belong to $M$.

If $S,U$ are initial sets, then $S+U$ is initial. For each pair $(s,u)$,
induct on the natural (Hessenberg) sum of their ordinal birthdays.
Every recursive addition option $s^L+u$, $s+u^L$, $s^R+u$ or $s+u^R$
has a smaller operand measure. It belongs to $S+U$ and, by induction,
so do all its prefixes. The preceding option-covering argument for the
defining cut of $s+u$ now puts all prefixes of $s+u$ in $S+U$.
The induction is on operand birthdays, not on the birthday of the sum.

Negation reverses all signs and preserves prefixes. Hence the additive
subgroup generated by any initial set containing zero is initial:
take finite sums of its union with its negative, and then their union.
It follows that, for an initial set-sized subgroup $M$ and a surreal
$b$ simplest between sets of options from $M$, the group $M+\Z b$
is initial. This directly supplies the additive adjunction needed here;
no general module-adjunction theorem is assumed.
If desired, when $b\notin M$, its small options extend to the partition
$M_{<b},M_{>b}$: $b$ realizes it, and any realization of the full partition
realizes the small cut, so $b$ is still simplest.

\subsection{Finite-support elements first}
First suppose $\Gamma=E$ is a set. Then $A$ is a set, since all its series
are encoded in the set of real-coefficient functions on subsets of $E$.
Well order $E$ by nondecreasing exponent birthday, breaking ties by choice.
All proper prefixes of an exponent precede it.
Starting with $\{0\}$, build initial subgroups by adjoining actual pure terms.

At exponent $y$, first adjoin $\omega^y$ if necessary.
All options in \eqref{leadercut} are in the previously generated group:
earlier exponent monomials give the left options, and
$\D\subseteq R_{y^R}$ gives the right ones.
The additive adjunction proves initiality and supplies all integer
multiples. Next adjoin the noninteger $r\omega^y$, $r\in R_y$,
in nondecreasing real birthday. Initiality of $R_y$ puts its real
predecessors in $R_y$; their multiples have already appeared
(integer multiples are already present).
For nondyadic $r$, or noninteger dyadic $r$ with no left prefix of $y$,
the pure cut uses just $r^L\omega^y$ and $r^R\omega^y$.
In the other noninteger dyadic case its options are
\[
 r^L\omega^y+k\omega^{y^L},\qquad
 r^R\omega^y-k\omega^{y^L},
\]
with $k\in\mathbb N_0$ and $y^L$ a left prefix of $y$.
They too are in the preceding group.
The pure-multiple cut and adjunction apply in every case.
At limits take unions. The resulting initial subgroup $F$ is exactly
the additive subgroup generated by all pure terms of $A$.
It contains every finite-support member of $A$, and only elements of $A$.

\subsection{Legal cuts for all actual limit-length supports}
Let $x=\sum_{\mu<\lambda}r_\mu\omega^{y_\mu}\in A$, with $\lambda$ an
infinite limit ordinal, and set
$s_\eta=\sum_{\mu<\eta}r_\mu\omega^{y_\mu}$.
For $\mu<\lambda$ put
\begin{equation}\label{unitcuts}
 l_\mu=s_{\mu+1}-\omega^{y_\mu},\qquad
 u_\mu=s_{\mu+1}+\omega^{y_\mu}.
\end{equation}
These are in $A$, by truncation closure and cross sectionality, and
$l_\mu<x<u_\mu$ by leading-coefficient comparison.
They use integer, not dyadic, perturbations.

They define $x$ as their simplest separator. Indeed, for any $\mu<\lambda$
and $k\ge1$, choose $\nu$ with $\mu<\nu<\lambda$. Normal-form comparison
gives
\begin{equation}\label{cofinal}
 s_{\mu+1}-2^{-k}\omega^{y_\mu}
 <l_\nu<x<u_\nu<
 s_{\mu+1}+2^{-k}\omega^{y_\mu}.
\end{equation}
Thus any separator of all the unit options is a separator of the standard
limit-series cut, whose simplest element is $x$.
As $x$ itself separates the unit options, it is their simplest separator.
We do not assert that every option is itself a prefix of $x$.

Well order all actual members of $A$ of infinite limit support length,
first by that length and then arbitrarily within each length.
Start from $F$ and adjoin them successively. At the stage for $x$ of
length $\lambda$, every member of $A$ of smaller support length is already
present. To see this, every ordinal $\rho$ is uniquely $\delta+j$,
where $j<\omega$ and $\delta$ is zero or a limit ordinal.
A series of length $\rho$ is its $\delta$-prefix, if nonempty, plus
$j$ pure terms. Its infinite-limit prefix was adjoined at a smaller
length, or it is already finite. Each option in \eqref{unitcuts} has
length at most $\mu+1<\lambda$, so belongs to the previous subgroup.
Adjunction applies. At limit stages take unions.
The same $\delta+j$ decomposition proves that the resulting initial
group is all of $A$. This includes supports cofinal in the entire
exponent enumeration; no bounded-stage union is used.

\subsection{Class exponents via actual set restrictions}
For an arbitrary exponent class and $x\in A$, its support $S$ is a set.
Let
\[
 E=\bigcup_{y\in S}\bigl(\{y\}\cup P(y)\bigr),
 \qquad A_E=\{b\in A:\supp(b)\subseteq E\}.
\]
For $x=0$, one may take $E=\varnothing$.
Replacement and union make $E$ a set, and initiality of $\Gamma$ makes
$E$ initial and contained in $\Gamma$. The group $A_E$ is a set, as
a subclass of the set of all real-coefficient series over $E$.
It inherits truncation closure, cross sectionality, the \emph{same}
pure coefficient groups at each $y\in E$, and the right-prefix condition.
The set proof makes $A_E$ initial. Since $x\in A_E$, every prefix of $x$
is in $A_E\subseteq A$. This proves Proposition~\ref{suff} for classes.
There is no class-sized cut or class enumeration.

For completeness, necessity uses the normal-form predecessor properties
of \cite[Proposition 3.2(i)--(ii), Lemma 3.1(i),(iv)]{ek}:
initiality gives truncation closure and initiality of the occurring
exponent class; subtraction of adjacent truncations isolates pure terms,
whose simpler leaders are present. It makes each $R_y$ initial and puts
$2^{-k}\omega^y$ in the group whenever $y\in R_s(x)$ and $\omega^x$ is
present. Additive closure gives $\D\subseteq R_y$.
These are the same explicit foundational imports as above.

\begin{thebibliography}{9}
\bibitem{thesis}
E.~Kaplan, \emph{Initial Embeddings in the Surreal Number Tree},
B.Sc.\ honors thesis, Ohio University, April 2015, \S7.1, p.~47.
\href{https://etd.ohiolink.edu/acprod/odb_etd/ws/send_file/send?accession=ouhonors1429615758&disposition=inline}{OhioLINK institutional full text}.
\bibitem{owr}
E.~Kaplan (joint work with P.~Ehrlich), Number systems with simplicity
hierarchies, in \emph{Surreal Numbers, Surreal Analysis, Hahn Fields and
Derivations}, Oberwolfach Reports 60/2016, Question 1, p.~3356
in the inspected EMS PDF (PDF p.~44).
\href{https://doi.org/10.4171/OWR/2016/60}{doi:10.4171/OWR/2016/60}.
\bibitem{ek}
P.~Ehrlich and E.~Kaplan, Number systems with simplicity hierarchies:
a generalization of Conway's theory of surreal numbers II,
\emph{J. Symbolic Logic} \textbf{83} (2018), 617--633.
\href{https://doi.org/10.1017/jsl.2017.9}{doi:10.1017/jsl.2017.9}.
Numbered statements here refer to the
\href{https://elliotakaplan.github.io/Number_systems_with_simplicity_hierarchies_II.pdf}{inspected 16-page author PDF};
the publisher body was not authenticated.
\href{https://arxiv.org/abs/1512.04001v1}{arXiv:1512.04001v1} is an earlier version.
\bibitem{adh}
M.~Aschenbrenner, L.~van den Dries and J.~van der Hoeven,
The surreal numbers as a universal $H$-field,
\emph{J. Eur. Math. Soc.} \textbf{21} (2019), 1179--1199.
Inspected \href{https://arxiv.org/abs/1512.02267v3}{arXiv:1512.02267v3}
(2016), Lemma 3.2, p.~11.
\bibitem{kks}
E.~Kaplan, L.~S.~Krapp and M.~Serra,
Decomposing the automorphism group of the surreal numbers,
\href{https://arxiv.org/abs/2509.22374v3}{arXiv:2509.22374v3}
(2026; first version 2025), Proposition 3.9 and Construction 3.10.
\bibitem{bvh}
V.~Bagayoko and J.~van der Hoeven, Surreal substructures,
\href{https://arxiv.org/abs/2305.02001}{arXiv:2305.02001} (2023).
Also inspected the author PDF dated 16 April 2022;
Theorem 6.7 and Proposition 7.1 are nearby general constructions.
\bibitem{handbook}
P.~Ehrlich, The absolute arithmetic continuum and its geometric
counterpart, in \emph{Handbook of the History and Philosophy of
Mathematical Practice} (2024), 1677--1718.
\href{https://doi.org/10.1007/978-3-031-40846-5_54}{doi:10.1007/978-3-031-40846-5\_54}.
Primary abstract and references only; full chapter unread.
\end{thebibliography}
\end{document}
